% GENERATED FILE. Edit canonical lessons, then run npm run generate:books.
% canonical-source-hash: fnv1a64:2845037105691c4c
% canonical-lessons: PA-C80-bura, PA-C80-lamma, PA-C80-garam, PA-C80-hauli, PA-C80-tez

\chapter{Bad, Long, Hot, Slowly, Fast}
\label{ch:pa-bad-long-hot-slowly-fast}
\hlchaptermodality{80}

\begin{chapteropening}
\textbf{By the end of this chapter:} \emph{I can say bad, long, tall, hot, slowly and fast.}

Complete the last lesson of chapter 80: I can say bad, long, tall, hot, slowly and fast.
\end{chapteropening}

\section[\texorpdfstring{\pa{ਬੁਰਾ} (burā)}{ਬੁਰਾ (burā)}]{\pa{ਬੁਰਾ} (burā) — bad}

\label{lesson:PA-C80-bura}



\begin{quote}
Before the new one: say the Punjabi for new, then the Punjabi for old.
\end{quote}

\subsection*{You'll want to know: \pa{ਬੁਰਾ}}
\textbf{\pa{ਬੁਰਾ}} — \emph{burā} — \textquotedblleft{}bad\textquotedblright{}.

\begin{cousinweb}
\textbf{\pa{ਬੁਰਾ}} (\emph{burā}) is a tadbhava, a word inherited and reshaped by sound change. It comes from \emph{*bura}, \textquotedblleft{}defective\textquotedblright{}, a form CDIAL reconstructs from its descendants (CDIAL 9289). CDIAL lists Punjabi \emph{burā}, \textquotedblleft{}wicked\textquotedblright{}.
\end{cousinweb}

\begin{grammarlens}[title={What you've built}]
One more word for this chapter.
\end{grammarlens}

\subsection*{Your turn}
\begin{itemize}
\raggedright
  \item \emph{Say it:} \emph{burā}
  \item \emph{Say it:} \emph{burā}, once more
  \item \emph{Say it:} say \emph{burā}
  \item \emph{Recall:} say the Punjabi for new, then the Punjabi for old, then say \emph{burā} again
\end{itemize}

\begin{tcolorbox}[breakable,colback=teal!4,colframe=teal!35!black,title={Before you move on}]
What does \pa{ਬੁਰਾ} mean? (\textquotedblleft{}Bad\textquotedblright{}.) Say it once more.
\end{tcolorbox}
\section[\texorpdfstring{\pa{ਲੰਮਾ} (lammā)}{ਲੰਮਾ (lammā)}]{\pa{ਲੰਮਾ} (lammā) — long, tall}

\label{lesson:PA-C80-lamma}



\begin{quote}
Before the new one: say the Punjabi for old, then the Punjabi for bad.
\end{quote}

\subsection*{You'll want to know: \pa{ਲੰਮਾ}}
\textbf{\pa{ਲੰਮਾ}} — \emph{lammā} — \textquotedblleft{}long, tall\textquotedblright{}.

\begin{cousinweb}
\textbf{\pa{ਲੰਮਾ}} (\emph{lammā}) is a tadbhava, a word inherited from Sanskrit and reshaped by sound change. It comes from \emph{lamba}, \textquotedblleft{}pendent\textquotedblright{} (CDIAL 10951). The Sanskrit word already meant \textquotedblleft{}long\textquotedblright{} in compounds.
\end{cousinweb}

\begin{grammarlens}[title={What you've built}]
One more word for this chapter.
\end{grammarlens}

\subsection*{Your turn}
\begin{itemize}
\raggedright
  \item \emph{Say it:} \emph{lammā}
  \item \emph{Say it:} \emph{lammā}, once more
  \item \emph{Say it:} say \emph{lammā}
  \item \emph{Recall:} say the Punjabi for old, then the Punjabi for bad, then say \emph{lammā} again
\end{itemize}

\begin{tcolorbox}[breakable,colback=teal!4,colframe=teal!35!black,title={Before you move on}]
What does \pa{ਲੰਮਾ} mean? (\textquotedblleft{}Long, tall\textquotedblright{}.) Say it once more.
\end{tcolorbox}
\section[\texorpdfstring{\pa{ਗਰਮ} (garam)}{ਗਰਮ (garam)}]{\pa{ਗਰਮ} (garam) — hot}

\label{lesson:PA-C80-garam}



\begin{quote}
Before the new one: say the Punjabi for bad, then the Punjabi for long.
\end{quote}

\subsection*{You'll want to know: \pa{ਗਰਮ}}
\textbf{\pa{ਗਰਮ}} — \emph{garam} — \textquotedblleft{}hot\textquotedblright{}.

\begin{cousinweb}
\textbf{\pa{ਗਰਮ}} (\emph{garam}) is a loan, usually traced to Persian \emph{garm}, \textquotedblleft{}warm, hot\textquotedblright{}. Punjabi breaks up the \emph{-rm} with a vowel.
\end{cousinweb}

\begin{grammarlens}[title={What you've built}]
One more word for this chapter.
\end{grammarlens}

\subsection*{Your turn}
\begin{itemize}
\raggedright
  \item \emph{Say it:} \emph{garam}
  \item \emph{Say it:} \emph{garam}, once more
  \item \emph{Say it:} say \emph{garam}
  \item \emph{Recall:} say the Punjabi for bad, then the Punjabi for long, then say \emph{garam} again
\end{itemize}

\begin{tcolorbox}[breakable,colback=teal!4,colframe=teal!35!black,title={Before you move on}]
What does \pa{ਗਰਮ} mean? (\textquotedblleft{}Hot\textquotedblright{}.) Say it once more.
\end{tcolorbox}
\section[\texorpdfstring{\pa{ਹੌਲੀ} (haulī)}{ਹੌਲੀ (haulī)}]{\pa{ਹੌਲੀ} (haulī) — slowly}

\label{lesson:PA-C80-hauli}



\begin{quote}
Before the new one: say the Punjabi for long, then the Punjabi for hot.
\end{quote}

\subsection*{You'll want to know: \pa{ਹੌਲੀ}}
\textbf{\pa{ਹੌਲੀ}} — \emph{haulī} — \textquotedblleft{}slowly\textquotedblright{}.

\begin{cousinweb}
\textbf{\pa{ਹੌਲੀ}} (\emph{haulī}) goes with \emph{haulā}, \textquotedblleft{}light, gentle\textquotedblright{}, which CDIAL lists under Sanskrit \emph{laghu}, \textquotedblleft{}light\textquotedblright{}, with its sounds swapped (CDIAL 10896).
\end{cousinweb}

\begin{grammarlens}[title={What you've built}]
One more word for this chapter.
\end{grammarlens}

\subsection*{Your turn}
\begin{itemize}
\raggedright
  \item \emph{Say it:} \emph{haulī}
  \item \emph{Say it:} \emph{haulī}, once more
  \item \emph{Say it:} say \emph{haulī}
  \item \emph{Recall:} say the Punjabi for long, then the Punjabi for hot, then say \emph{haulī} again
\end{itemize}

\begin{tcolorbox}[breakable,colback=teal!4,colframe=teal!35!black,title={Before you move on}]
What does \pa{ਹੌਲੀ} mean? (\textquotedblleft{}Slowly\textquotedblright{}.) Say it once more.
\end{tcolorbox}
\section[\texorpdfstring{\pa{ਤੇਜ਼} (tez)}{ਤੇਜ਼ (tez)}]{\pa{ਤੇਜ਼} (tez) — fast}

\label{lesson:PA-C80-tez}



\begin{quote}
Before the new one: say the Punjabi for hot, then the Punjabi for slowly.
\end{quote}

\subsection*{You'll want to know: \pa{ਤੇਜ਼}}
\textbf{\pa{ਤੇਜ਼}} — \emph{tez} — \textquotedblleft{}fast\textquotedblright{}.

\begin{cousinweb}
\textbf{\pa{ਤੇਜ਼}} (\emph{tez}) is a loan, usually traced to Persian \emph{tēz}, \textquotedblleft{}sharp, swift\textquotedblright{}.
\end{cousinweb}

\begin{grammarlens}[title={What you've built}]
One more word for this chapter.
\end{grammarlens}

\subsection*{Your turn}
\begin{itemize}
\raggedright
  \item \emph{Say it:} \emph{tez}
  \item \emph{Say it:} \emph{tez}, once more
  \item \emph{Say it:} say \emph{tez}
  \item \emph{Recall:} say the Punjabi for hot, then the Punjabi for slowly, then say \emph{tez} again
\end{itemize}

\begin{tcolorbox}[breakable,colback=teal!4,colframe=teal!35!black,title={Before you move on}]
What does \pa{ਤੇਜ਼} mean? (\textquotedblleft{}Fast\textquotedblright{}.) Say it once more.
\end{tcolorbox}
